\documentclass{article}
\usepackage{graphicx} 
\usepackage{cite}
\usepackage{amsmath,amssymb,amsfonts}
\usepackage{algorithmic}
\usepackage{textcomp}
\usepackage{xcolor}
\usepackage{booktabs}
\usepackage{arydshln}

%Add more packages if needed

\title{Laboratory \#1}
\author{Authors}
\date{April 2026}

\begin{document}

\maketitle

\section{Introduction}

To cite a paper you can use~\cite{AgudoPAMI2022,YanICCV2015}.


\section{Methods}


The matrix can be factorized into two terms including intrinsic and extrinsic parameters, respectively, as:
$
\mathbf{P}_i=	\begin{bmatrix}
	\alpha_x & s & \beta_x\\
	0 & \alpha_y & \beta_y\\
	0&0&1
\end{bmatrix}
\begin{bmatrix}
	r_{11} & r_{12} & r_{13} & t_{x}\\
	r_{21} & r_{22} & r_{23} & t_{y}\\
	r_{31} & r_{32} & r_{33} & t_{z}
\end{bmatrix},
$
with $(\alpha_x,\alpha_y)$ the focal length, $(\beta_x,\beta_y)$ the principal point and $s$ the aspect ratio.

or
\begin{equation}
\mathbf{P}_i=	\begin{bmatrix}
	\alpha_x & s & \beta_x\\
	0 & \alpha_y & \beta_y\\
	0&0&1
\end{bmatrix}
\begin{bmatrix}
	r_{11} & r_{12} & r_{13} & t_{x}\\
	r_{21} & r_{22} & r_{23} & t_{y}\\
	r_{31} & r_{32} & r_{33} & t_{z}
\end{bmatrix},
\end{equation}

\newpage

\section{Experiments}

An example of figure:

\begin{figure}[t!]
\centering
\resizebox{8.5cm}{!} {
\begin{tabular}{@{}ccc@{}}
%\hspace{-0.2cm}
  \includegraphics[clip, angle=0,width=1.0\linewidth]{Figures/result1.png}&
\hspace{-0.15cm}
\includegraphics[clip, angle=0,width=1.0\linewidth]{Figures/result2.png}&
\hspace{-0.15cm}
  \includegraphics[clip, angle=0,width=1.0\linewidth]{Figures/result3.png}
\end{tabular}}
\vspace{-0.2cm}
   \caption{\textbf{Qualitative matching evaluation}. Some instances of }
\label{fig:experiment1}
%\vspace{-0.2cm}
\end{figure}




\section{Conclusion}



Please, consider the maximum number of pages per laboratory. 




\bibliographystyle{IEEEbib}
\bibliography{references}

\end{document}
